\chapter{相关技术与研究现状}
\label{chap:background}

本论文研究小模型如何依据知识库的执行状态继续完成问答，以及专业资料中的证据如何进入可核验的回答。二者分别涉及训练分布、交互决策和来源表达。本章选择与这些问题直接相关的原始文献，说明已有技术提供了什么条件、本文沿用了什么，以及哪些结论仍需要本地实验支持。文献中的模型规模、任务和指标不一致，本文不将其报告分数与后续实验直接排名。

\section{知识图谱问答与显式程序表示}

知识图谱问答将自然语言需求转化为知识库上的实体定位、关系查询或属性操作。显式程序在模型与知识库之间提供一个可检查的接口：模型预测操作及参数，执行器按照确定规则计算结果。相比只输出答案，程序还允许保存中间结果、检查依赖和回放失败，但程序合法不代表它完整表达了问题意图。

Cao 等提出的 KQA Pro 包含约十二万道复杂知识库问题，为问题提供 KoPL 程序和 SPARQL 查询。KoPL 用组合操作表达推理，覆盖关系、字面值及限定信息；原文附录给出了属性查询、限定条件查询、集合操作和比较等函数\cite{cao2022kqapro}。这使研究者可以区分最终答案、程序结构和具体操作能力，为本文构造可执行训练轨迹提供基础。

本文使用 KQA Pro 的数据和执行任务，但后续五千训练题、五百开发题与两千项目留出题是项目内部划分，不是原论文的完整标准评测。雷达应用只实现较窄的属性查询接口，也不能因工具名称与 KoPL 相近就称其具有完整的多跳、集合或数值推理能力。来源读法的条件保留与公开任务的限定查询具有联系，二者仍需分别验证。

\section{小语言模型与参数高效适配}

Qwen2.5 技术报告介绍了多种规模的开放权重模型及其指令版本，其中包含三十亿参数量级模型，并说明了预训练、监督微调与后训练过程\cite{qwen2024qwen25}。本文选择 Qwen2.5-3B-Instruct 作为固定模型基础，是为了在有限计算资源下进行配对实验。模型已有的语言和代码能力并不保证它会使用本项目的函数签名，基础模型报告的能力也不能代替专用任务评价。

LoRA 通过冻结原权重、训练低秩增量来完成适配\cite{hu2022lora}。对权重矩阵 $W_0\in\mathbb{R}^{d\times k}$，可将更新写为
\begin{equation}
 W=W_0+\frac{\alpha}{r}UV,\qquad
 U\in\mathbb{R}^{d\times r},\quad V\in\mathbb{R}^{r\times k},
\end{equation}
其中 $r$ 为秩，$\alpha$ 控制缩放。训练的主要变化量由低秩矩阵承担，因而不必为每项适配更新全部原始权重。

本文沿用该技术实现训练，不将调整秩、学习率或微批量本身视为方法贡献。研究变量是模型见到哪些执行状态、哪些输出位置接受监督。比较这些变量时还须区分监督 token 数、完整输入长度和实际计算成本：前者相同只限制目标信息量，不自动意味着训练时间或推理开销相等。

\section{工具交互与执行反馈}

ReAct 将语言推理与环境动作交替组织，使模型能够在生成过程中获取外部信息并调整后续步骤。原文既研究少样本提示，也报告了用成功轨迹微调较小模型的实验\cite{yao2023react}。因此，“使用工具”和“训练工具使用能力”都已有明确研究基础，不能仅凭增加交互轮次就宣称提出新范式。

本文的分步问答同样依赖动作和观察的交替，但主要输出是受约束的知识操作及结束动作，不把自由文本反思作为必需步骤。执行反馈包括查询结果、句柄和失败信息。成功反馈仍可能来自错误实体、遗漏条件或无关查询；被执行器拒绝的操作则只能说明接口不合法。两种情况对学习的意义不同，不能统一称为语法错误。

由此形成两个待检验问题：模型能否在实际自由运行状态后继续完成任务，以及改善是否依赖某一类反馈内容。前者可比较最终成功率和完成率；后者需要反馈干预或组成部分消融。观察到无效调用减少，尚不足以证明模型理解了报错语义，也不能把可读轨迹直接当作内部推理机制的证据。

\section{序列决策中的分布偏移与恢复监督}

在顺序决策中，当前动作会改变后续观察。只在参考轨迹上训练时，模型见到的状态分布与自由运行形成的分布可能不同。Ross 等的 DAgger 在学习策略访问的状态上获取专家动作，并迭代聚合数据，用监督学习改善当前策略；其分析在相应假设下将这一过程联系到无遗憾在线学习\cite{ross2011dagger}。该工作提示，训练样本的状态来源与动作标签同样重要。

本文与这一动机相关，但不实施完整 DAgger：只从一个固定初始模型采集训练自由轨迹，不反复采样新策略，也没有能在任意偏离状态给出最佳动作的通用专家。可用监督来自原有参考程序。构造器保留共同成功前缀，插入真实分歧动作及观察，再对参考后缀进行句柄重映射和执行核验；无法继续执行的样本被排除。因此其适用范围受可恢复状态覆盖限制，不能直接继承 DAgger 的理论保证。

这里的首次分歧仅相对于一条参考程序定义。不同程序可能正确回答同一问题，故分歧不等于已被标注的语义错误。本文将分歧作为上下文而非模仿目标，监督恢复后的操作；与之配对的普通续训使用相同问题和参考后缀。该设计用于检验整体恢复构造的增量效果，尚不能分离额外状态、依赖句柄变化和监督位置各自的贡献。

\section{反思、自我修订与强化学习的区别}

Reflexion 将任务反馈转成语言反思，并将其保存在记忆中影响后续尝试，其核心方案不通过更新模型权重实现这种改进\cite{shinn2023reflexion}。Self-Refine 则让同一个语言模型对初始输出给出反馈，再反复修订；该方案不要求额外监督训练\cite{madaan2023selfrefine}。它们说明推理时利用反馈有多种形式，但语言形式的自我评价仍须与外部执行结果区分。

本文的主要恢复方案在训练阶段更新 LoRA 参数，反馈由实际执行器产生，目标后缀经确定性重放验证。它不训练模型生成反思文本，也不在评价时要求额外的自我批评回合。因此应将其定位为恢复状态下的监督学习构造，而不是把 Reflexion 改名后复现，或声称已经实现通用自我纠错。

强化学习还涉及另一种优化方式。DeepSeekMath 中的 GRPO 对同一问题采样多条输出，用组内奖励形成相对优势，避免另外拟合一个价值模型\cite{shao2024deepseekmath}。本文只将其作为监督训练后的有限扩展对照。原论文的数学任务收益不能推出知识库任务必有收益；本地短程扩展的负面或不确定结果，也不能推出该算法普遍无效。采用反馈不等于采用强化学习，普通交叉熵恢复监督与奖励优化应分别命名和评价。

\section{检索增强、证据组织与来源引用}

Lewis 等提出的 RAG 将参数化生成器与外部检索记忆结合，原方法使用稠密检索和序列到序列模型，并在训练中组合检索结果与生成概率\cite{lewis2020rag}。本文领域基线采用固定 BM25 页块检索和同一回答模型，属于检索增强的应用路线，并非复现原论文的联合训练架构。保留这一参照，是为了判断结构化整理的额外成本是否带来可观察的回答收益。

雷达资料常在表格中同时出现多个型号、部件、量值形式和操作条件。来源表示若丢掉这些关系，即使数值没有抄错，答案也可能错配。本文将完整平铺记录与条件绑定组织作主要表示对照，要求二者可无损还原同一关系元组序列。原文页含有额外未结构化内容，故只作系统参照。删除条件后再比较准确率改变了信息量，不能用来单独证明条件组织算法有效。

ALCE 将回答正确性与引用质量分开，并分别评价引用对陈述的支持覆盖及无关引用情况\cite{gao2023alce}。这一视角说明，存在引用标识、引用属于已提供资料、资料支持正文，是不同层次。本文借鉴分层评价思想，但采用题级正文、实际证据支持和引用支持的独立判断，不称为原 ALCE 指标的直接复现，也没有把每个冗余但合法引用单独计为失败。

引用格式还可能影响联合指标。若模型显式引用记录 ID，而该 ID 唯一映射到实际提供的来源页，可以用统一规则检查这种标识差异的影响；但不能依据参考答案补引用、猜测缺失来源或改写正文。第六章将此类解析保留为事后敏感性分析，与冻结的原始分数并列。它检验工程解析对指标的影响，不构成新的事实推理能力。

\section{研究定位与方法比较}

表\ref{tab:background-methods}按反馈来源、是否更新参数和本文的关系比较相关方法。表中说明技术区别，不比较他文与本文的准确率；本文没有复现其中所有方法，因此不能根据概念比较宣称总体优于它们。

\begin{longtable}{@{}p{0.17\linewidth}p{0.22\linewidth}p{0.23\linewidth}p{0.28\linewidth}@{}}
\caption{相关方法与本文的比较维度}\label{tab:background-methods}\\
\toprule
方法或路线 & 可用反馈与数据 & 主要作用位置 & 本文关系及限制\\
\midrule
\endfirsthead
\toprule
方法或路线 & 可用反馈与数据 & 主要作用位置 & 本文关系及限制\\
\midrule
\endhead
KoPL 程序问答\cite{cao2022kqapro} & 问题、参考程序与执行结果 & 将语义需求映射到可执行操作 & 提供任务和监督基础；雷达接口只实现其中较窄能力\\
ReAct\cite{yao2023react} & 交替的语言推理、动作与观察 & 提示交互；原文也研究轨迹微调 & 沿用交互思路，本文不要求生成显式反思\\
DAgger\cite{ross2011dagger} & 当前策略访问状态上的专家动作 & 迭代聚合数据并训练策略 & 动机相近；本文只有固定采集与可重放后缀，不继承理论保证\\
Reflexion / Self-Refine\cite{shinn2023reflexion,madaan2023selfrefine} & 环境反馈或模型自评与修订 & 核心方案在推理上下文中利用反馈 & 本文主要更新参数，未建立反思文本机制\\
GRPO\cite{shao2024deepseekmath} & 同题多条采样及相对奖励 & 以策略优化更新参数 & 只作有限扩展；不将其称为恢复监督的必要组成\\
本文恢复配对监督 & 实际分歧状态与可执行参考后缀 & 后缀位置的监督微调 & 匹配问题、次序及监督 token；输入长度仍不同\\
本文条件绑定表示 & 同一检索页对应的完整语义记录 & 固定生成器前的确定性组织 & 无新训练；与引用解析共同考察，不宣称普遍领域推理提升\\
\bottomrule
\end{longtable}

论文的主方法问题是：在既定公开任务和匹配监督预算下，引入真实分歧状态的恢复后缀训练，是否优于普通续训。领域问题则是：当可信来源、主体和条件明确后，如何让查询结果进入有据回答，以及同样信息的组织方式是否影响条件绑定与引用。两个问题通过数据流程衔接，但拥有不同模型配置、评价材料和证据强度，不能把领域表示结果追溯归因于公开恢复训练。

评价可信度也属于方法设计。可执行程序能够核对动作与结果，却不能证明资料解读无误；两个 AI 审阅者一致仍不等于人工金标；从一份资料派生的多题也不等于多份独立来源。后续章节保留这些限制，并将预先规定的评价、开发诊断和事后分析分别报告，避免仅用单一总分解释全部效果。

\section{本章小结}

已有研究分别提供了显式程序问答、参数高效适配、交互式工具使用、偏离状态监督和引用评价的基础。本文的改进空间是在可回放知识库任务中构造并控制恢复监督，以及在雷达资料中保持证据关系和引用可核验性。下一章首先明确数据与来源边界，第四、五章检验恢复训练，随后在第六章评价领域接口和证据表示；各阶段的结论由对应实验支撑，不以技术名称替代证据。
